% =================================================
% Ученический рабочий лист
% =================================================

\worksheettitle
  {Логарифм: в какую степень?}
  {Рабочий лист: открываем новое понятие через степени}

\studentnameblock

\begin{checkpointbox}[Входная диагностика: 2 минуты]
  Заполните пропуски без калькулятора.
  \begin{enumerate}[label=\textbf{\arabic*.},itemsep=0.48em]
    \item \(2^{\mathgap[6mm]}=32\).
    \item \(10^{\mathgap[6mm]}=0{,}01\).
    \item \(3^{-2}=\answerblank[22mm]\).
    \item \(25^{\frac12}=\answerblank[18mm]\).
  \end{enumerate}
\end{checkpointbox}



\section{Показатель можно подобрать не всегда}

Сначала решим два уравнения:
\[
  2^x=8,
  \qquad
  2^x=\frac14.
\]

\begin{fillbox}[Подберите показатели]
  \[
8=2^{\mathgap[6mm]}
    \quad\Longrightarrow\quad
    x=\answerblank[14mm],
  \]
  \[
    \frac14=2^{\mathgap[6mm]}
    \quad\Longrightarrow\quad
    x=\answerblank[14mm].
  \]
\end{fillbox}

Теперь рассмотрим уравнение
\[
  2^x=5.
\]

\begin{thinkbox}[Оцените корень до введения нового обозначения]
  Используйте соседние степени двойки:
  \[
    2^{\mathgap[6mm]}=4<5<8=2^{\mathgap[6mm]}.
  \]
  Это значит, что
  \[
    2^2<2^x<2^3.
  \]
  Функция \(y=2^t\) строго возрастает, поэтому
  \[
    \answerblank[14mm]<x<\answerblank[14mm].
  \]
\end{thinkbox}

\logmotivationworksheetfigure

\begin{checkpointbox}[Что показывает рисунок]
  \begin{enumerate}
    \item Сколько точек пересечения имеют графики? \answerblank[18mm]
    \item Функция \(y=2^x\) строго возрастает. Объясните, почему поэтому
      уравнение \(2^x=5\) имеет только один корень.
      \writinglines{1}
    \item Можно ли по рисунку определить точное значение корня?
      \answerblank[25mm]
  \end{enumerate}
\end{checkpointbox}

По графику можно установить существование и примерное положение корня,
но нельзя определить его точное значение. Для точной записи вводят новое
обозначение
\[
  x_0=\log_2 5.
\]

\begin{reflectionbox}[Проговорите смысл]
  Дополните фразу:

  \(\log_2 5\) --- это такой \answerblank[34mm] степени числа 2,
  при котором получается число \answerblank[14mm].
\end{reflectionbox}

\Needspace{13\baselineskip}
\section{Определение логарифма}

\begin{fillbox}[Восстановите определение]
  Логарифмом числа \(b\) по основанию \(a\) называют
  \answerblank[42mm] степени, в которую нужно возвести число
  \answerblank[14mm], чтобы получить число \answerblank[14mm].
\end{fillbox}

\logequivalenceworksheetfigure

\begin{checkpointbox}[Главная формула]
  Впишите одну и ту же букву в оба пропуска и завершите равносильность:
  \[
    \boxed{
      \log_a b=\answerblank[12mm]
      \quad\Longleftrightarrow\quad
      a^{\answerfraction[12mm]}=b
    }.
  \]
\end{checkpointbox}

\begin{fillbox}[Роли чисел в записи]
  В равенстве \(\log_a b=c\):
  \begin{itemize}
    \item \(a\) --- \answerblank[35mm] степени;
    \item \(b\) --- число, которое нужно \answerblank[38mm];
    \item \(c\) --- искомый \answerblank[35mm] степени.
  \end{itemize}
\end{fillbox}

\begin{notebox}[Условия записи]
  В вещественных числах логарифм рассматривают при
  \[
    a>0,\qquad a\ne1,\qquad b>0.
  \]
  При вычислении логарифмов далее мы будем использовать только корректные
  основания и аргументы. В следующем задании специально приведены также
  недопустимые записи. Причины ограничений будут изучены отдельно.
\end{notebox}

\begin{checkpointbox}[Допустима ли запись?]
  Отметьте записи, которые имеют смысл в вещественных числах, и для
  каждой недопустимой записи укажите нарушенное условие:
  \[
    \log_2 8,
    \qquad
    \log_{-2}8,
    \qquad
    \log_2(-8),
    \qquad
    \log_1 5.
  \]
  \writinglines{2}
\end{checkpointbox}

\section{Переводим запись в обе стороны}

\begin{fillbox}[Из степени в логарифм]
  По образцу заполните правый столбец.

  \begin{center}
    \renewcommand{\arraystretch}{1.48}
    \begin{tabularx}{\textwidth}{
      >{\centering\arraybackslash}p{0.37\textwidth}
      >{\centering\arraybackslash}p{0.10\textwidth}
      >{\centering\arraybackslash}X}
      \toprule
      \rowcolor{TableHeader}
      Степенная запись & & Логарифмическая запись \\
      \midrule
      \(2^3=8\) & \(\Longleftrightarrow\) & \(\log_2 8=3\) \\
      \(5^2=25\) & \(\Longleftrightarrow\) & \(\log_{\mathgap[12mm]}\mathgap[12mm]=\mathgap[12mm]\) \\
      \(10^{-3}=0{,}001\) & \(\Longleftrightarrow\) & \(\log_{\mathgap[12mm]}\mathgap[12mm]=\mathgap[12mm]\) \\
      \(16^{\frac12}=4\) & \(\Longleftrightarrow\) & \(\log_{\mathgap[12mm]}\mathgap[12mm]=\mathgap[12mm]\) \\
      \(\left(\frac13\right)^{-2}=9\) & \(\Longleftrightarrow\) & \(\log_{\mathgap[12mm]}\mathgap[12mm]=\mathgap[12mm]\) \\
      \bottomrule
    \end{tabularx}
  \end{center}
\end{fillbox}

\begin{fillbox}[Из логарифма в степень]
  \begin{center}
    \renewcommand{\arraystretch}{1.75}
    \begin{tabular}{@{}r@{\hspace{0.55em}}r@{\hspace{0.45em}}c@{\hspace{0.45em}}l@{\hspace{1.35em}}c@{\hspace{1.35em}}c@{\hspace{0.75em}}c@{\hspace{0.75em}}c@{}}
      \textbf{1.} & \(\log_4 64\) & \(=\) & \(3\) &
        \(\Longleftrightarrow\) &
        \(\answerblank[14mm]^{\answerfraction[11mm]}\) & \(=\) &
        \(\answerblank[18mm].\) \\
      \textbf{2.} & \(\log_7 1\) & \(=\) & \(0\) &
        \(\Longleftrightarrow\) &
        \(\answerblank[14mm]^{\answerfraction[11mm]}\) & \(=\) &
        \(\answerblank[18mm].\) \\
      \textbf{3.} & \(\log_{25}5\) & \(=\) & \(\displaystyle\frac12\) &
        \(\Longleftrightarrow\) &
        \(\answerblank[14mm]^{\answerfraction[11mm]}\) & \(=\) &
        \(\answerblank[18mm].\) \\
    \end{tabular}
  \end{center}
\end{fillbox}

\begin{reflectionbox}[Проверка перевода]
  В равенстве \(a^c=b\) число \(a\) становится
  \answerblank[28mm] логарифма, число \(c\) —
  \answerblank[28mm] логарифма, а число \(b\) —
  \answerblank[28mm] логарифма.
\end{reflectionbox}

\Needspace{13\baselineskip}
\section{Вычисляем логарифмы}

\begin{methodintro}[Один вопрос вместо новой формулы]
  Чтобы вычислить \(\log_a b\), спросите:
  \begin{center}
    \textbf{«В какую степень нужно возвести \(a\), чтобы получить \(b\)?»}
  \end{center}
\end{methodintro}

\begin{fillbox}[Заполните таблицу рассуждений]
  \begin{center}
    \small
    \renewcommand{\arraystretch}{1.55}
    \begin{tabularx}{\textwidth}{
      >{\centering\arraybackslash}p{0.19\textwidth}
      >{\centering\arraybackslash}p{0.29\textwidth}
      >{\centering\arraybackslash}p{0.25\textwidth}
      >{\centering\arraybackslash}X}
      \toprule
      \rowcolor{TableHeader}
      Выражение & Вопрос & Представление числа & Ответ \\
      \midrule
      \(\log_2 32\) & \(2^{?}=32\) & \(32=2^5\) & \(5\) \\
      \(\log_5 1\) & \(5^{?}=1\) & \(1=5^{\answerfraction[8mm]}\) & \answerblank[10mm] \\
      \(\log_3\dfrac1{27}\) & \(3^{?}=\dfrac1{27}\) & \(\dfrac1{27}=3^{\answerfraction[8mm]}\) & \answerblank[10mm] \\
      \(\log_9 3\) & \(9^{?}=3\) & \(3=9^{\answerfraction[8mm]}\) & \answerblank[10mm] \\
      \(\log_{\frac12}8\) & \(\left(\dfrac12\right)^{?}=8\) & \(8=\left(\dfrac12\right)^{\answerfraction[8mm]}\) & \answerblank[10mm] \\
      \bottomrule
    \end{tabularx}
  \end{center}
\end{fillbox}

\begin{checkpointbox}[Что заметили]
  Значение логарифма может быть:
  \choicebox\ положительным;
  \quad\choicebox\ нулём;
  \quad\choicebox\ отрицательным;
  \quad\choicebox\ дробным.
\end{checkpointbox}
\newpage
\begin{practicebox}[Самостоятельно: без готовых шагов]
  Вычислите
  \[
    \log_{\frac13}27.
  \]
  Запишите вопрос, полное рассуждение и обязательную степенную проверку
  ответа.
  \workarea{28mm}
\end{practicebox}

\section{Решаем простейшие уравнения}

\begin{fillbox}[Уравнение 1: неизвестен аргумент логарифма]
  \[
    \log_2 x=6
    \quad\Longleftrightarrow\quad
    2^{\answerfraction[12mm]}=x
    \quad\Longrightarrow\quad
    x=\answerblank[20mm].
  \]
\end{fillbox}

\begin{fillbox}[Уравнение 2: неизвестное содержится в аргументе логарифма]
  \[
    \log_5(x+1)=2
    \quad\Longleftrightarrow\quad
    x+1=\answerblank[20mm]
    \quad\Longrightarrow\quad
    x=\answerblank[20mm].
  \]
\end{fillbox}

\begin{fillbox}[Уравнение 3: неизвестен показатель]
  Значение показателя в уравнении \(3^x=10\) не удаётся непосредственно
  подобрать по известным степеням числа 3. По определению логарифма его
  точный корень равен
  \[
    x=\answerblank[36mm].
  \]
\end{fillbox}

\begin{notebox}[Точное и приближённое значение]
  Запись \(x=\log_3 10\) задаёт точное значение корня. Десятичное
  приближение можно получить с помощью научного калькулятора или
  математической программы:
  \[
    \log_3 10\approx2{,}096,
    \qquad
    \text{до сотых: }\log_3 10\approx2{,}10.
  \]
\end{notebox}

\begin{questionsbox}[Итоговое задание]
  Выполните самостоятельно, не возвращаясь к образцам.
  \begin{enumerate}[itemsep=0.45em]
    \item Объясните словами смысл записи \(\log_2 7\).
    \item Переведите равенство \(3^{-2}=\dfrac19\) в логарифмическую форму.
    \item Проверьте равенство \(\log_4 64=3\) степенным равенством.
  \end{enumerate}
\end{questionsbox}

\begin{checkpointbox}[Обязательная часть завершена]
  Вы должны уметь:
  \begin{itemize}
    \item объяснить смысл записи \(\log_a b\);
    \item переводить степенную запись в логарифмическую и обратно;
    \item вычислять простые логарифмы через представление числа степенью;
    \item проверять значение логарифма соответствующим степенным
      равенством \(a^c=b\);
    \item распознавать допустимые основания и аргументы;
    \item записывать точный корень уравнения \(a^x=b\).
  \end{itemize}
\end{checkpointbox}

\newpage
\section*{Расширение}

\section{Найдите и исправьте ошибки}

\begin{practicebox}[Ошибка 1]
  Ученик написал:
  \[
    \log_4 64=16,
    \qquad \text{потому что }4\cdot16=64.
  \]
  Какой вопрос нужно было задать? Запишите верное решение.
  \workarea{30mm}
\end{practicebox}

\begin{practicebox}[Ошибка 2]
  Из равенства \(2^5=32\) ученик получил \(\log_{32}2=5\).
  Объясните, какие роли чисел перепутаны, и запишите верное равенство.
  \workarea{28mm}
\end{practicebox}

\begin{practicebox}[Ошибка 3]
  Ученик утверждает: «Логарифм не может быть отрицательным, потому что
  степень --- это количество множителей». Опровергните утверждение одним
  примером.
  \workarea{24mm}
\end{practicebox}

\newpage
\section*{Основные тождества}

\section{Получаем первые общие равенства}

\begin{fillbox}[Логарифм единицы]
  Для любого допустимого основания \(a\):
  \[
    a^{\answerfraction[12mm]}=1
    \quad\Longrightarrow\quad
    \boxed{\log_a1=\answerblank[12mm]}.
  \]
\end{fillbox}

\begin{fillbox}[Логарифм основания]
  \[
    a^{\answerfraction[12mm]}=a
    \quad\Longrightarrow\quad
    \boxed{\log_a a=\answerblank[12mm]}.
  \]
\end{fillbox}



\section{Восстановите полное рассуждение}

\begin{practicebox}[Задание 1]
  Вычислите \(\log_{\frac14}64\). Решение должно содержать:
  \begin{enumerate}
    \item вопрос о показателе степени;
    \item представление числа 64 степенью основания \(\dfrac14\);
    \item проверку полученного ответа.
  \end{enumerate}
  \workarea{48mm}
\end{practicebox}

\begin{practicebox}[Задание 2]
  Решите уравнение
  \[
    \log_3(2x-1)=4.
  \]
  \workarea{40mm}
\end{practicebox}

\begin{practicebox}[Задание 3]
  Запишите точный корень уравнения
  \[
    7^x=20
  \]
  и укажите два соседних целых числа, между которыми он расположен.
  \workarea{38mm}
\end{practicebox}

\begin{reflectionbox}[Моё определение]
  Не используя учебник, завершите предложение:

  «Логарифм числа \(b\) по основанию \(a\) --- это
  \answerblank[98mm]».
\end{reflectionbox}

\begin{questionsbox}[Перед завершением работы]
  \begin{enumerate}
    \item Что нужно проверить, чтобы убедиться в равенстве \(\log_a b=c\)?
    \item Как из равенства \(a^c=b\) получить логарифмическую запись?
    \item Почему \(\log_5 1=0\)?
    \item Как записать точный корень уравнения \(11^x=6\)?
  \end{enumerate}
\end{questionsbox}
